% Source: Wald GR, physical PDF pp. 321--325 (printed pp. 308--312).
% Coverage: Section 12.2, General Properties of Black Holes; theorems 12.2.1, 12.2.5--12.2.6 and propositions 12.2.2--12.2.4.
% Figures/tables/footnotes: footnote 2; no figures or tables.
% Status: source-reviewed and compile-clean.
% Uncertainties: none.

\subsection{General Properties of Black Holes}\label{sec:12.2}
\setcounter{theorem}{0}

In this section we shall use the global methods developed in chapters 8 and 9 to
establish some general results in the theory of black holes. Most of these results are
due to Hawking. Our discussion will differ slightly from that of Hawking and Ellis
(1973) in that we shall not impose topological restrictions on the Cauchy surfaces for
\(M\cap\widetilde V\), and some of our definitions (such as that of an outer trapped
surface) are not identical to theirs.

First, we note that for any asymptotically flat spacetime \((M,g_{ab})\) with
associated unphysical spacetime \((\overline M,\widetilde g_{ab})\), if
\(q\in\mathcal I^+\) and \(p\in M\cap J^-(q)\), then any point \(r\in\mathcal I^+\)
lying beyond \(q\) on the future directed null geodesic generator of \(\mathcal I^+\)
passing through \(q\), satisfies \(p\in I^-(r)\). (This follows from the corollary to
theorem~8.1.2, since \(p\) and \(r\) can be joined by a causal curve which is not an
unbroken null geodesic.) Hence, we have
\(M\cap J^-(\mathcal I^+)=M\cap I^-(\mathcal I^+)\), so
\(J^-(\mathcal I^+)\) is open in \(M\). Thus, the black hole region,
\(B=M-J^-(\mathcal I^+)\), is closed in \(M\). In particular this means that the
event horizon \(H\) is contained in \(B\).

We define, now, the notion of a black hole at a given instant of time. Let
\((M,g_{ab})\) be a strongly asymptotically predictable spacetime, with globally
hyperbolic region \(\widetilde V\supseteq\overline M\cap J^-(\mathcal I^+)\) in the
unphysical spacetime. Let \(B=M-J^-(\mathcal I^+)\) be the black hole region of the
spacetime. If \(\Sigma\) is a Cauchy surface for \(\widetilde V\), we shall refer to
\(\Sigma\cap B\) as the total black hole region at time \(\Sigma\). Each connected
component, \(\mathcal B\), of \(\Sigma\cap B\) will be called a \emph{black hole at
time \(\Sigma\)}. The number of black holes present in \((M,g_{ab})\) may vary with
\lq\lq time\rq\rq\ (i.e., choice of Cauchy surface \(\Sigma\)), since new black holes
may form and black holes present at one time may later coalesce. However, our first
theorem may be interpreted as stating that a black hole may never disappear nor may
it \lq\lq bifurcate,\rq\rq\ i.e., split into more than one black hole at a later time.

\begin{theorem}\label{thm:12.2.1}
Let \((M,g_{ab})\) be a strongly asymptotically predictable spacetime and let
\(\Sigma_1\) and \(\Sigma_2\) be Cauchy surfaces for \(\widetilde V\) with
\(\Sigma_2\subseteq I^+(\Sigma_1)\). Let \(\mathcal B_1\) be a nonempty connected
component of \(B\cap\Sigma_1\), i.e., \(\mathcal B_1\) is a black hole at time
\(\Sigma_1\). Then \(J^+(\mathcal B_1)\cap\Sigma_2\ne\emptyset\) and is contained
within a single connected component of \(B\cap\Sigma_2\).
\end{theorem}
\begin{proof}
That \(J^+(\mathcal B_1)\cap\Sigma_2\ne\emptyset\) follows immediately from the
fact that \(\Sigma_2\) is a Cauchy surface lying to the future of \(\Sigma_1\).
Clearly \(J^+(\mathcal B_1)\subseteq B\), so \(J^+(\mathcal B_1)\cap\Sigma_2\) is
contained within \(B\cap\Sigma_2\). Therefore, to complete the proof of the
theorem, it suffices to show that \(J^+(\mathcal B_1)\cap\Sigma_2\) is connected.
Suppose this were not the case. Then we could find disjoint open sets
\(O,O'\subseteq\Sigma_2\) with
\(O\cap J^+(\mathcal B_1)\ne\emptyset\), \(O'\cap J^+(\mathcal B_1)\ne\emptyset\),
and \(O\cup O'\supseteq J^+(\mathcal B_1)\cap\Sigma_2\). Then we would have
\(\mathcal B_1\cap I^-(O)\ne\emptyset\), \(\mathcal B_1\cap I^-(O')\ne\emptyset\),
and \(\mathcal B_1\subseteq I^-(O)\cup I^-(O')\). However, no point, \(p\), of
\(\mathcal B_1\) can lie in both \(I^-(O)\) and \(I^-(O')\), for then we could
divide the future directed timelike geodesics from \(p\) according to whether they
intersected \(\Sigma_2\) in \(O\) or \(O'\) and thereby divide the timelike vectors at
\(p\) into nonempty, disjoint open sets. (This would contradict the fact that the
interior of the future light cone of \(T_pM\) is connected.) Thus,
\(I^-(O)\cap\Sigma_1\) and \(I^-(O')\cap\Sigma_1\) must be disjoint open subsets
of \(\Sigma_1\), each of which intersects \(\mathcal B_1\) and whose union contains
\(\mathcal B_1\). However, this contradicts the hypothesis that \(\mathcal B_1\) is
connected.
\end{proof}

One of the difficulties commonly encountered in working with the definition of a
black hole is that one needs to know the entire future development of a spacetime in
order to determine if a black hole is present. However, it may happen that one is
given only initial data \((\Sigma,h_{ab},K_{ab})\) for a spacetime and cannot
explicitly solve the evolution equations. In that case, the general definition of a
black hole is of little use for directly determining whether \(B\ne\emptyset\) in the
spacetime determined by these data and, if \(B\ne\emptyset\), where on \(\Sigma\)
the black hole region \(\mathcal B=B\cap\Sigma\) may lie. Therefore, it is useful to
develop criteria for the existence and location of a black hole without requiring
knowledge of the global time development of the spacetime. The next three results
give us such criteria.

In chapter 9, we proved that if a trapped surface is present in a spacetime,
appropriate energy conditions are satisfied by matter, and a number of further
hypotheses hold, then a singularity must occur (see theorems~9.5.3 and~9.5.4). Thus,
trapped surfaces are associated with spacetime singularities. Our first result giving
criteria for existence of a black hole states that in a strongly asymptotically
predictable spacetime it is not only true that no singularities can be \lq\lq seen\rq\rq\
from null infinity, but it also is true that no trapped surfaces are visible from null
infinity. In other words, all trapped surfaces must be entirely contained within black
holes.

\begin{proposition}\label{prop:12.2.2}
Let \((M,g_{ab})\) be a strongly asymptotically predictable spacetime for which
\(R_{ab}k^ak^b\ge0\) for all null \(k^a\), as will be the case if Einstein's equation
holds with the weak or strong energy condition satisfied by matter. Suppose \(M\)
contains a trapped surface, \(T\). Then \(T\subseteq B\), where \(B\) denotes the
black hole region of the spacetime.
\end{proposition}
\begin{proof}
Suppose \(T\) were not entirely contained within \(B\). Then in the unphysical
spacetime, we would have \(J^+(T)\cap\mathcal I^+\ne\emptyset\). However, spatial
infinity, \(\ii^0\), is not to the causal future of any point in \(M\), so clearly
\(\ii^0\notin J^+(T)\). Furthermore, since the region \(\widetilde V\) of the
unphysical spacetime is globally hyperbolic and \(T\) is compact, it follows from
theorem~8.3.11 that \(J^+(T)\) is closed in \(\widetilde V\). Therefore there is an
open neighborhood of \(\ii^0\) which fails to intersect \(J^+(T)\), and, thus, an open
region of \(\mathcal I^+\) does not intersect \(J^+(T)\). Since \(\mathcal I^+\) is
connected, this implies that there exists a point \(q\in\mathcal I^+\) such that
\(q\in\partial J^+(T)\). Hence, according to theorem~9.3.11, in the unphysical
spacetime there is a null geodesic \(\gamma\) from \(p\in T\) to \(q\) which is
orthogonal to \(T\) and has no conjugate point between \(T\) and \(q\). With respect
to the physical metric, \(g_{ab}\), \(\gamma\) also is a null geodesic (see appendix
D) orthogonal to \(T\) with no conjugate point, but now \(\gamma\) is future
complete. However, this is impossible, because according to theorem~9.3.6, in the
physical spacetime \(\gamma\) must have a conjugate point within affine parameter
\(2/|\theta_0|\) from \(p\), where \(\theta_0<0\) is the expansion at \(p\) of the
orthogonal null geodesic congruence from \(T\) to which \(\gamma\) belongs.
\end{proof}

In fact, by a somewhat different argument, we can slightly generalize
proposition~12.2.2 to apply to a \emph{marginally trapped surface}, i.e., a compact,
spacelike two-dimensional submanifold for which the expansion of both families of
orthogonal geodesics is required only to be nonpositive, \(\theta\le0\), rather than
strictly negative, \(\theta<0\). We have

\begin{proposition}\label{prop:12.2.3}
Let \(T\) be a marginally trapped surface in a strongly asymptotically predictable
spacetime for which \(R_{ab}k^ak^b\ge0\) for all null \(k^a\). Then \(T\subseteq B\).
\end{proposition}
\begin{proof}
As in the proof of the previous proposition, we know from theorem~9.3.11 that
\(\partial J^+(T)\) is generated by the null geodesics orthogonal to \(T\). The
expansion, \(\theta\), of these null geodesics in the physical spacetime is
nonpositive initially. Therefore, we have \(\theta\le0\) everywhere on
\(\partial J^+(T)\), since \(\dd\theta/\dd\lambda\le0\) according to
equation~(9.2.32) and \(\theta\) cannot become positive on \(\partial J^+(T)\) by
passing through \(-\infty\) because that would imply existence of a conjugate point.
Suppose, now, that \(T\cap J^-(\mathcal I^+)\ne\emptyset\). Then, as in the proof
of the previous proposition, in the unphysical spacetime there would exist a point
\(q\in\mathcal I^+\) with \(q\in\partial J^+(T)\). Furthermore, all the points
lying to the causal future of \(q\) along the null geodesic generator of
\(\mathcal I^+\) passing through \(q\) must lie in \(I^+(T)\), since they can be
joined to \(T\) by a causal curve which is not an unbroken null geodesic. Since
\(I^+(T)\) is open, all generators of \(\mathcal I^+\) sufficiently near the one
passing through \(q\) must enter \(I^+(T)=\operatorname{int}(J^+(T))\). On the
other hand, all generators of \(\mathcal I^+\) have past endpoints at \(\ii^0\) and
thus leave \(J^+(T)=\overline{J^+(T)}\). Therefore, not only \(q\) but also an
entire local cross section, \(\mathcal S\), of \(\mathcal I^+\) must lie on
\(\partial J^+(T)\). Furthermore, the null geodesic generators of
\(\partial J^+(T)\) must strike \(\mathcal S\) orthogonally in order that there not
exist any timelike curves from \(T\) to \(\mathcal S\). However, in the physical
spacetime, the expansion of the null geodesic congruence orthogonal to any cross
section of \(\mathcal I^+\) is positive near \(\mathcal I^+\). This contradicts the
previous result that the generators of \(\partial J^+(T)\) have nonpositive expansion
everywhere.
\end{proof}

Note that the only properties of \(T\) needed in the proofs of
propositions~12.2.2 and~12.2.3 are that \(J^+(T)\) is closed, that
\(\ii^0\notin J^+(T)\), and that the expansion of the null geodesic generators of
\(\partial J^+(T)\) is initially nonpositive. Another important example of a set
having these properties is the following. Let \(\Sigma\) be any asymptotically flat
Cauchy surface for \(\widetilde V\), so that \(\Sigma\) passes through \(\ii^0\) and is
spacelike there. Let \(C\subseteq\Sigma\cap M\) be a closed subset of \(\Sigma\)
which forms a three-dimensional manifold with boundary (see appendix B) and
suppose the two-dimensional boundary \(S=\partial C\) of \(C\) has the property
that the expansion, \(\theta\), of the outgoing family of null geodesics orthogonal to
\(S\) is everywhere nonpositive, \(\theta\le0\). (Here, the \emph{outgoing family}
is defined to be the family of null geodesics orthogonal to \(S\) satisfying
\(k^aN_a\ge0\), where \(k^a\) denotes a geodesic tangent and \(N^a\) is the normal
to \(S\) in \(\Sigma\) which points outward from \(C\).) A surface \(S\) satisfying
these properties is called an \emph{outer marginally trapped surface}, and \(C\) is
called a \emph{trapped region}. (Note that \(C\) need not be connected or compact.
Note also that a trapped surface need not be outer trapped since it need not be the
boundary of a three-dimensional volume.) Then \(C\) can be seen to satisfy the above
desired properties as follows: By problem 8 of chapter 8, \(J^+(C)\) is closed.
Clearly, we have \(\ii^0\notin J^+(C)\). Finally, each of the null geodesic generators
of \(\partial J^+(C)\subseteq J^+(C)\) must have a past endpoint on \(C\) but
cannot meet \(C\) in \(\operatorname{int}(C)\), cannot meet \(S=\partial C\)
nonorthogonally, and cannot be an ingoing geodesic at \(S\), since in any of these
cases it would enter \(I^+(C)\). Thus, \(\theta\le0\) initially for the null geodesic
generators of \(\partial J^+(C)\). Hence, a repetition of the proof of
proposition~12.2.3 yields the following result.

\begin{proposition}\label{prop:12.2.4}
Let \((M,g_{ab})\) be a strongly asymptotically predictable spacetime for which
\(R_{ab}k^ak^b\ge0\) for all null \(k^a\). Let \(\Sigma\) be an asymptotically flat
Cauchy surface for \(\widetilde V\), and let \(C\subseteq\Sigma\) be a trapped
region. Then \(C\subseteq B\cap\Sigma\).
\end{proposition}

We define the \emph{total trapped region}, \(\mathcal T\), of a Cauchy surface,
\(\Sigma\), to be the closure of the union of all trapped regions, \(C\), on
\(\Sigma\). We call the boundary, \(\mathcal A=\partial\mathcal T\), of
\(\mathcal T\) the \emph{apparent horizon} on \(\Sigma\). It follows immediately
from proposition~12.2.4 that in a strongly asymptotically predictable spacetime with
\(R_{ab}k^ak^b\ge0\) for all null \(k^a\), we have
\(\mathcal T\subseteq B\cap\Sigma\), so the apparent horizon always lies inside of
(or coincides with) the true event horizon, \(H\cap\Sigma\), on \(\Sigma\). The
apparent horizon, \(\mathcal A\), satisfies the following property.

\begin{theorem}\label{thm:12.2.5}
If the total trapped region, \(\mathcal T\), on a Cauchy surface \(\Sigma\) has the
structure of a manifold with boundary, then the apparent horizon, \(\mathcal A\), is
an outer marginally trapped surface with vanishing expansion, \(\theta=0\).
\end{theorem}
\begin{proof}
Clearly \(\mathcal T\) is closed and we also have\footnote{To see this we note that in
an asymptotically Euclidean coordinate system on \(\Sigma\) (see problem 2 of
chapter 11) there exists a radius \(R\) such that all coordinate spheres with
\(r>R\) have everywhere positive expansion of the outgoing null geodesics. Hence no
outer trapped surface, \(S\), can enter the region \(r>R\), since the expansion of the
outgoing null geodesics from \(S\) would have to be at least as great as that from the
coordinate sphere at the point where \(r\) takes its maximum value on \(S\). Thus,
\(\mathcal T\) cannot enter the region \(r>R\).} \(\ii^0\notin\mathcal T\), so
\(\mathcal T\) satisfies the first two requirements for being a trapped region. To show
that \(\theta=0\) on \(\mathcal A\), we note that if \(\theta>0\) at
\(p\in\mathcal A\), then we could find a neighborhood, \(U\), of \(p\) such that
every surface, \(S\), passing through \(U\) with \(\theta\le0\) everywhere on \(S\)
must leave \(\mathcal T\). Hence, we could not have outer trapped surfaces passing
arbitrarily close to \(p\) but staying within \(\mathcal T\) as required by the fact
that \(p\) is on the boundary of \(\mathcal T\). On the other hand, given that
\(\theta\le0\) everywhere on \(\mathcal A\), if we had \(\theta<0\) at
\(q\in\mathcal A\), we could deform \(\mathcal A\) outward in a neighborhood of
\(q\), preserving \(\theta\le0\) everywhere. In this way, we would produce a
trapped region \(C\) larger than \(\mathcal T\), which is impossible.
\end{proof}

The final result we shall prove in this section concerns the evolution of the event
horizon. Since the horizon, \(H\), is the boundary of the past of \(\mathcal I^+\), by
theorem~8.1.3 it is an achronal, three-dimensional, embedded \(C^0\) (in fact,
\(C^{1-}\)) submanifold. Furthermore, by theorem~8.1.6, \(H\) is generated by
future inextendible null geodesics, since no null geodesic generator of \(H\) can
have a future endpoint on \(\mathcal I^+\). Thus, if the intersection
\(\mathcal H=H\cap\Sigma\) of the horizon with a spacelike Cauchy surface,
\(\Sigma\), for \(\widetilde V\) is nonempty, it comprises a two-dimensional
submanifold of \(\Sigma\). The next theorem states that the area of \(\mathcal H\)
never decreases with time. As we shall see in Sections~12.5 and~14.4, this result
underlies what appears to be a profound relationship between black holes,
thermodynamics, and quantum physics.

\begin{theorem}[black hole area theorem; Hawking 1971]\label{thm:12.2.6}
Let \((M,g_{ab})\) be a strongly asymptotically predictable spacetime satisfying
\(R_{ab}k^ak^b\ge0\) for all null \(k^a\). Let \(\Sigma_1\) and \(\Sigma_2\) be
spacelike Cauchy surfaces for the globally hyperbolic region \(\widetilde V\) with
\(\Sigma_2\subseteq I^+(\Sigma_1)\) and let
\(\mathcal H_1=H\cap\Sigma_1\), \(\mathcal H_2=H\cap\Sigma_2\), where \(H\)
denotes the event horizon, i.e., the boundary of the black hole region of
\((M,g_{ab})\). Then the area of \(\mathcal H_2\) is greater than or equal to the
area of \(\mathcal H_1\).
\end{theorem}
\begin{proof}
We establish, first, that the expansion, \(\theta\), of the null geodesic generators of
\(H\) is everywhere nonnegative, \(\theta\ge0\). Suppose \(\theta<0\) at
\(p\in H\). Let \(\Sigma\) be a spacelike Cauchy surface for \(\widetilde V\) passing
through \(p\) and consider the two-surface \(\mathcal H=H\cap\Sigma\). Since
\(\theta<0\) at \(p\), we can deform \(\mathcal H\) outward in a neighborhood of
\(p\) to obtain a surface \(\mathcal H'\) on \(\Sigma\) which enters
\(J^-(\mathcal I^+)\) and has \(\theta<0\) everywhere in
\(J^-(\mathcal I^+)\). However, by the same type of argument as given in
proposition~12.2.2, this leads to a contradiction as follows. Let
\(K\subsetneq\Sigma\) be the closed region lying between \(\mathcal H\) and
\(\mathcal H'\), and let \(q\in\mathcal I^+\) with
\(q\in\partial J^+(K)\). Then the null geodesic generator of
\(\partial J^+(K)\) on which \(q\) lies must meet \(\mathcal H'\) orthogonally.
However, this is impossible, since \(\theta<0\) on \(\mathcal H'\), and thus this
generator will have a conjugate point before reaching \(q\). Thus, \(\theta\ge0\)
everywhere on \(H\).

Now, as mentioned above, each \(p\in\mathcal H_1\) lies on a future inextendible
null geodesic, \(\gamma\), contained in \(H\). Since \(\Sigma_2\) is a Cauchy
surface, \(\gamma\) must intersect \(\Sigma_2\) at a point \(q\in\mathcal H_2\).
Thus, we obtain a natural map from \(\mathcal H_1\) into (a portion of)
\(\mathcal H_2\). Since \(\theta\ge0\), the area of the portion of \(\mathcal H_2\)
given by the image of \(\mathcal H_1\) under this map must be at least as large as
the area of \(\mathcal H_1\). In addition, since the map need not be onto---e.g., new
black holes may form between \(\Sigma_1\) and \(\Sigma_2\)---the area of
\(\mathcal H_2\) may be even larger. Thus, the area of \(\mathcal H_2\) cannot be
smaller than that of \(\mathcal H_1\).
\end{proof}
